%% Por qué un pipe: de la tráquea (vía de conducción, imagen del thesis structure) al pipe periódico de la DNS.
%% Imágenes: figs/vias_conductoras.png (fondo transparente) y figs/pipe_contorno.png (contorno de u, research plan).
\begin{tikzpicture}[font=\small]
  % --- izquierda: vía aérea de conducción ---
  \node[anchor=south west,inner sep=0pt] (L) at (0,0)
    {\includegraphics[height=4.1cm]{figs/vias_conductoras.png}};
  % tráquea resaltada (x ≈ 53 %, y ≈ 82–100 % de la imagen)
  \coordinate (T) at ($(L.south west)!0.53!(L.south east) + (0,0.91*4.1cm)$);
  % toda la tráquea, de la boca a la carina (y ≈ 60–100 % de la imagen)
  \path (L.south west) ++(0.42*4.02cm,0.59*4.1cm) coordinate (T0) (L.south west) ++(0.63*4.02cm,4.1cm+0.08cm) coordinate (T1);
  \draw[ingEst,line width=1.2pt,rounded corners=3pt,fill=ingEst,fill opacity=0.10] (T0) rectangle (T1);
  \coordinate (T) at ($(T0)!0.5!(T1)$);
  \node[anchor=north,align=center,font=\footnotesize] at ($(L.south)+(0,-0.12)$)
    {tráquea adulta\\ {\color{tenue}$D \approx 18$ mm · 60 L/min · $\mathrm{Re}_\tau \approx 160$}};

  % --- flecha ---
  \draw[-{Stealth[length=7pt]},tenue,line width=1pt] ($(T1|-T)+(0.15,0.3)$) to[out=10,in=170]
    node[above,pos=0.5,font=\scriptsize,text=tenue]{hay DNS} (6.95,3.75);

  % --- derecha arriba: contorno de velocidad de una simulación ---
  \begin{scope}[shift={(7.2,3.75)}]
    \node[anchor=west,inner sep=0pt] (P) at (0,0) {\includegraphics[width=4.9cm,height=0.9cm]{figs/pipe_contorno.png}};
    \draw[-{Stealth[length=5pt]},tenue,dashed] ($(P.north east)+(-0.1,0.1)$) to[out=155,in=25]
      node[above,font=\scriptsize,text=tenue]{periódico en $z$} ($(P.north west)+(0.1,0.1)$);
  \end{scope}
  % --- derecha abajo: el caso (tubo con perfil turbulento, como en el thesis structure) ---
  \begin{scope}[shift={(7.2,1.9)}]
    \def\Lp{4.9}\def\Rp{0.55}\def\ex{0.22}
    \shade[top color=black!20, bottom color=black!20, middle color=black!3]
      (0,-\Rp) -- (\Lp,-\Rp) arc[start angle=-90,end angle=90,x radius=\ex,y radius=\Rp] -- (0,\Rp) -- cycle;
    \fill[fondo] (0,0) ellipse[x radius=\ex,y radius=\Rp];
    \foreach \y/\l in {-0.38/0.65, -0.19/1.1, 0/1.3, 0.19/1.1, 0.38/0.65}
      \draw[-{Stealth[length=4pt]},ingEsq,line width=0.7pt] (0,\y) -- ++(\l,0);
    \draw[ingEsq,line width=0.7pt] (0,-0.46) -- (0,0.46);
    \draw[tinta,line width=0.6pt] (0,-\Rp) -- (\Lp,-\Rp) arc[start angle=-90,end angle=90,x radius=\ex,y radius=\Rp] -- (0,\Rp);
    \draw[tinta,line width=0.6pt] (0,0) ellipse[x radius=\ex,y radius=\Rp];
    \draw[tinta,line width=0.6pt] (\Lp,0) -- (\Lp,\Rp) node[midway,right=3pt,font=\footnotesize]{$R = 9$ mm};
    \node[font=\footnotesize,text=tinta] at (3.3,0) {$f_z = 2u_\tau^2/R$};
    \draw[{Stealth[length=4pt]}-{Stealth[length=4pt]},tenue] (0,-\Rp-0.22) -- (\Lp,-\Rp-0.22)
      node[midway,below,font=\footnotesize,text=tinta]{$L = 15R$};
  \end{scope}
  \node[anchor=north,align=center,font=\footnotesize] at (9.65,0.55)
    {$\mathrm{Re}_\tau = 180$ · $\mathrm{Re}_b \approx 5300$ · $Q \approx 69$ L/min\\
     {\color{tenue}DNS: El Khoury et al.\ 2013}};
\end{tikzpicture}
